\section{The exact original-field witness and potential unlinking}
\label{sec:cubic-original-field-witness}

Use the actual Type I point, original field \(K\), and
representations \(\rho,\nu=\rho c_p^{-1}\) of AC-G111.
Retain all its original support and published arithmetic
inputs. Let \(\Pi\) be the original-field regular algebraic
conjugate self-dual cusp constructed in AC-G111-descent.
The local input below supplies the literal stronger
\texttt{WitA} condition, without applying the incompatible
cyclotomic \texttt{LevelAt} wrapper to the cubic type.

\paragraph{Exact published local input.}
We use Barnet-Lamb--Gee--Geraghty--Taylor,
\emph{Local-global compatibility for \(l=p\), II},
Theorem~1.1 and Corollary~3.2, in the author version
\url{https://www.ma.imperial.ac.uk/~tsg/Index_files/localglobalII.pdf}.
For a regular algebraic essentially conjugate self-dual
cuspidal representation of Shin-regular weight, the
conclusion is full Frobenius-semisimplified local--global
compatibility at coefficient primes, with the twist
\(|\det|^{(1-n)/2}\). Rank three is automatically
Shin-regular by the definition: there is no additional
gap condition on the highest weight in odd rank.
There is no Iwahori-fixed-vector or residual-irreducibility
hypothesis in this statement. Full compatibility away
from the coefficient prime is the already retained
published \texttt{IsGalRepOf.away} input.

We also use the all-finite-place temperedness explicitly
invoked in the complete proof of Proposition~1.1(2) of
the same authors' \emph{Local-global compatibility for
\(l=p\), I}, printed p.~5,
\url{https://www.ma.imperial.ac.uk/~tsg/Index_files/localglobalI.pdf}.
That proof cites Caraiani's Theorem~1.2 for temperedness
and Taylor--Yoshida for purity. We retain those published
theorems. The complete printed proofs of that proposition
and Corollary~1.3, and of II, Theorems~2.1 and~3.1 and
Corollary~3.2, were read. The cited Shimura-variety,
eigenvariety, and period-comparison theorems inside those
proofs are not being re-proved here.

Write \(|\cdot|_v\) for the norm character of the local
Weil group, taking geometric Frobenius to \(q_v^{-1}\).
The cyclotomic character has this Weil--Deligne character.
The native rank-three attachment normalization therefore
gives, at every finite place,
\[
 \begin{split}
 \operatorname{WD}(\nu|_{G_{K_v}})^{F\text{-ss}}
   &\simeq \operatorname{rec}(\Pi_v)\otimes|\cdot|_v^{-1},\\
 \operatorname{WD}(\rho|_{G_{K_v}})^{F\text{-ss}}
   &\simeq \operatorname{rec}(\Pi_v).
 \end{split}
 \tag{G111W.1}
\]
The second equality cancels the twist from
\(\nu=\rho c_p^{-1}\). It does not identify the
untwisted and normalized Galois representations.

\begin{theorem}[AC-G111-witness: the actual stronger witness]
\label{AC-G111-witness}
For every actual Type I point under these hypotheses,
including \(e=1\), the following statements hold.
\begin{enumerate}
\item At every finite place, the right side of (G111W.1)
is a direct sum of three unitary one-dimensional Weil
characters, with monodromy zero. In the literal source
notation, \(\operatorname{WinAt}(\Pi,0,v)\) holds for
every finite \(v\).
\item At a coefficient-prime place, take the actual full
ordered flag characters \(\kappa_0,\kappa_1,\kappa_2\)
of \(\rho\). For all \(i,j\) and all integers \(m\ge1\),
\[
                  \kappa_i^m\ne(\kappa_j c_p)^m.
 \tag{G111W.2}
\]
Thus the exact source \texttt{PotUnl} holds for both
\(\rho\) and \(\nu\). It continues to hold after every
finite local extension. In particular \texttt{Unl1} holds.
\item The literal \(\operatorname{WitA}(\iota,\nu)\)
condition holds. Together with AC-G111-descent this gives
\(\operatorname{IsAutCsdA}(\iota,\nu)\), the stronger
automorphic field in \texttt{Adm3Alt}. The older away
condition \texttt{Unlinked} also holds for \(\rho\)
and \(\nu\).
\end{enumerate}
No \texttt{LevelAt}, \texttt{FlagAt}, \texttt{Adm3U\_s},
or uniform counting conclusion is included in this theorem.
\end{theorem}
\begin{proof}
\emph{1. The actual local Weil parameters are toral.}
Away from \(p\), the modified killing condition makes
inertia factor through the finite abelian prime-to-\(p\)
image of the residual tuple. Since this image has order
invertible in the coefficient integers, its character
idempotents are integral. Their ranks are unchanged on
reduction. Uniqueness of prime-to-\(p\) root lifts shows
that its three characters, with multiplicities, are
precisely the Teichmuller lifts of the restrictions of
the three global residual characters.

Every element of the decomposition group preserves these
inertia isotypic spaces: the inertia characters just
described extend to global characters, so they are
unchanged by conjugation. On a repeated isotypic space
inertia is scalar. Diagonalizing the semisimple Frobenius
on each such space therefore splits the Frobenius
semisimplification into one-dimensional Weil characters.
The inertia action is finite, so the monodromy operator
of the Galois Weil--Deligne representation is zero.
This also covers a repeated residual inertia character;
no assumption of three distinct characters at every
away place is imposed.

At \(p\), the actual AC-G65 component is potentially
crystalline, so again the monodromy is zero. Its full
ordered local flag is already proved for every point of
that component. The category of potentially crystalline
representations is closed under subquotients, and the
Weil--Deligne functor is exact on it. Applied to this
flag, it shows that the Frobenius-semisimplified Weil
representation is the direct sum of the three
one-dimensional parameters of \(\kappa_i\).
This argument uses the actual flag and component, not
merely a prescribed inertial semisimplification.

Equation (G111W.1) now gives the same description for
\(\operatorname{rec}(\Pi_v)\) at every finite place.
Since \(\Pi_v\) is tempered, a zero-monodromy parameter
which is a sum of one-dimensional Weil representations
has unitary characters: their restrictions to inertia
are finite and their values on geometric Frobenius
have complex absolute value one. Equivalently the
associated local principal-series cuspidal support
consists of three unitary characters of \(K_v^\times\).
This is the standard normalized local Langlands
description of a tempered parameter with monodromy
zero; there is no nontrivial special-representation
block left when monodromy is zero.

Consequently the source \texttt{TorAt} predicate holds,
and every rank-one character in every cuspidal support
has value of absolute value one on a uniformizer.
Since \(q_v>1\),
\[
                      q_v^{-1/2}<1<q_v^{1/2}.
 \tag{G111W.3}
\]
These are exactly all inequalities and the torality
clause in \texttt{WinAt} at \(t=0\). This proves the
first assertion, including its quantifier over every
cuspidal support.

\emph{2. Exclude every potential cyclotomic link.}
At \(p\), denote the one-dimensional Weil character
of \(\kappa_i\) by \(\xi_i\). Step~1 gives
\(|\iota\xi_i(\operatorname{Frob}_v)|=1\).
If (G111W.2) failed, exact tensor functoriality of the
Weil--Deligne construction for one-dimensional de Rham
characters would give
\[
                    \xi_i^m=\xi_j^m|\cdot|_v^m.
 \tag{G111W.4}
\]
On geometric Frobenius the complex absolute values of
the two sides are respectively \(1\) and \(q_v^{-m}\),
which differ. This is a contradiction. In particular
the proof does not require a gap \(e\ge2\); the
potentially problematic adjacent weights when \(e=1\)
are handled by this same absolute-value argument.

Restriction to a finite local extension restricts the
unitary Weil characters and the norm character. The
same argument uses its residue cardinality and proves
(G111W.2) there. Alternatively, if a ratio became
finite order on an open subgroup, a further positive
power of it would be trivial on the original local
group, contrary to what was just proved. Multiplying
every \(\kappa_i\) by the same scalar character leaves
the comparison in (G111W.2) unchanged. It therefore
also applies to the flag of \(\nu\). Conjugating a
place or its decomposition group transports the full
flag and this argument, so the assertion has the
literal source quantifier over all choices of a place
above the given prime.

\emph{3. Check both exact witness clauses.}
The first clause of \texttt{WitA} is the implication
\texttt{Unl1} \(\Rightarrow\) \texttt{PotUnl} at
coefficient primes. Step~2 proves its conclusion
directly for \(\nu\). For the second clause, let
\(\Pi'\) be any regular algebraic cuspidal
representation attached to \(\nu\) in the native
\texttt{AttC} sense. Its almost-everywhere Satake system
is the same as that of \(\Pi\). Strong multiplicity
one gives \(\Pi'\simeq\Pi\). Step~1 consequently
gives \(\operatorname{WinAt}(\Pi',0,v)\) at every
finite place, hence in particular at every place in
the disjunction specified by \texttt{WitA}.
This is the required single real number \(t\), namely
zero. The representation is already continuous, so
the continuity implication in \texttt{IsAutCsdA} adds
no premise.

For clarity, the older \texttt{Unlinked} predicate can
also be read directly. At an away place take an
arithmetic Frobenius lift. Its \(\rho\)-eigenvalues
have absolute value one by Step~1; inversion exchanges
arithmetic and geometric Frobenius and preserves this
absolute value. Finite inertia only changes them by
unitary factors. No two eigenvalues can have positive
powers differing by the corresponding positive power
of \(q_v>1\). For \(\nu\), all absolute values are
multiplied by the same scalar, so the same conclusion
holds. Thus one may take the Frobenius-power index
\(1\) in the literal \texttt{UnlinkAt} definition.
\end{proof}

\paragraph{Source dictionary and exact boundary.}
The literal definitions used here are \texttt{WinAt}
in \path{LeviSmoothLocalization/WinUpDown.lean},
\texttt{Unl1}, \texttt{PotUnl}, and \texttt{WitA}
in \path{AutomorphyLifting/WeakOfAtt.lean}, and
\texttt{IsAutCsdA} in \path{AutomorphyLifting/IsAutCsdA.lean}.
Their containing files were read in full. The older
\texttt{NoLink}/\texttt{UnlinkAt}/\texttt{Unlinked}
definitions occur in \path{IsAutomorphicOverImagQuad.lean}.
The source \texttt{Th.Unl} additionally packages a
\texttt{FlagAt} condition with cyclotomic finite
characters. It is not identified with the coefficient-
independent \texttt{Unl1}/\texttt{PotUnl} conclusions
above. The order-three inertia obstruction to that
cyclotomic wrapper remains exactly as stated in
AC-G111. The separate bounded local-type counting
argument must retain the actual cubic inertia.
